% =========================================================
\section{Schnellentscheidungshilfen (Prüfungs-Kompass)}
% =========================================================

Diese Section dient als schnelle Entscheidungslogik für typische Prüfungsfragen.
Sie ersetzt keine Herleitungen, sondern liefert direkte \textbf{Wenn–Dann}-Aussagen.

% ---------------------------------------------------------
\subsection{Allgemeine Grundlagen (R, L, C)}

\subsubsection{Grundrelationen}

\[
u_R = R i,\qquad u_L = L\frac{di}{dt},\qquad i_C = C\frac{du_C}{dt}
\]

\begin{itemize}
\item Wenn $i$ in einer Spule sprunghaft ändern soll $\Rightarrow$ $u_L$ wird sehr gross $\Rightarrow$ Freilaufpfad nötig.
\item Wenn $u_C$ sprunghaft ändern soll $\Rightarrow$ $i_C$ wird sehr gross $\Rightarrow$ Einschaltstromspitze.
\item Induktivität: \textbf{Strom ist quasi-kontinuierlich}.
\item Kapazität: \textbf{Spannung ist quasi-kontinuierlich}.
\end{itemize}

\subsubsection{Zeitkonstante und Einschwingen}

\[
\tau_{RL}=\frac{L}{R},\qquad \tau_{RC}=RC
\]

\begin{itemize}
\item Wenn $\tau$ gross $\Rightarrow$ langsame Änderung, starke Glättung.
\item Wenn $\tau$ klein $\Rightarrow$ schnelle Änderung, wenig Glättung.
\item Nach $t\approx 5\tau$ ist das System praktisch im Endwert.
\end{itemize}

% ---------------------------------------------------------
\subsection{Mittelwert, Effektivwert, Leistung}

\subsubsection{Schnellregeln}

\[
\overline{x}=\frac{1}{T}\int_0^T x(t)\,dt,\qquad
X_{\mathrm{RMS}}=\sqrt{\frac{1}{T}\int_0^T x^2(t)\,dt}
\]

\begin{itemize}
\item Wenn Signal symmetrisch um 0 (ungerade) $\Rightarrow$ Mittelwert $=0$.
\item Wenn Signal immer positiv (Gleichrichtung) $\Rightarrow$ Mittelwert $>0$.
\item RMS reagiert stark auf Spitzenwerte $\Rightarrow$ Impulsströme erhöhen RMS stark.
\end{itemize}

\subsubsection{Leistung}

\[
p(t)=u(t)i(t),\quad P=\overline{p(t)}
\]

\begin{itemize}
\item Ohmsche Last: $P=\frac{U_{\mathrm{RMS}}^2}{R}=R I_{\mathrm{RMS}}^2$.
\item Induktiv/Kapazitiv: Momentanleistung kann positiv/negativ sein $\Rightarrow$ Blindleistung möglich.
\end{itemize}

\subsubsection{Nichtsinusförmig}

\[
U_{\mathrm{RMS}}^2=\sum_{n\ge 1} U_{n,\mathrm{RMS}}^2
\]

\begin{itemize}
\item Wenn THD hoch $\Rightarrow$ Verzerrung hoch $\Rightarrow$ Filter/Glättung sinnvoll.
\item Wenn Strom verzerrt $\Rightarrow$ Leistungsfaktor sinkt auch bei $\cos\varphi_1\approx 1$.
\end{itemize}

% ---------------------------------------------------------
\subsection{Leistungshalbleiter (Diode, Thyristor, MOSFET, IGBT)}

\subsubsection{Grundlogik}

\begin{itemize}
\item Diode: leitet automatisch, keine Steuerung.
\item Thyristor: zündbar, nicht abschaltbar (nur über Stromnull oder Kommutierung).
\item MOSFET: voll steuerbar, schnell, gut für hohe $f_s$.
\item IGBT: voll steuerbar, gut bei hohen Leistungen, mittlere $f_s$.
\end{itemize}

\subsubsection{Verluste}

\[
P_{cond}=U_0 I_{\mathrm{avg}} + r I_{\mathrm{RMS}}^2,\qquad
P_{sw}=f_s(E_{on}+E_{off})
\]

\begin{itemize}
\item Wenn $f_s\uparrow$ $\Rightarrow$ Schaltverluste $\uparrow$.
\item Wenn $I_{\mathrm{RMS}}\uparrow$ $\Rightarrow$ Leitverluste $\uparrow$.
\item MOSFET: Verluste dominieren über $r_{DS,on}I^2$.
\item IGBT: Verluste dominieren über $U_{CE0}I$ bei mittleren Strömen.
\item Diode: hohe $Q_{rr}$ $\Rightarrow$ hohe Schaltverluste in schnellen Umrichtern.
\end{itemize}

\subsubsection{Thermik}

\[
T_j=T_{amb}+R_{th,tot}P_V,\qquad
R_{th,tot}=R_{th,jc}+R_{th,ck}+R_{th,ka}
\]

\begin{itemize}
\item Wenn $P_V\uparrow$ $\Rightarrow$ $T_j\uparrow$.
\item Wenn $R_{th}\downarrow$ (besserer Kühlkörper) $\Rightarrow$ $T_j\downarrow$.
\item Grenzcheck immer: $T_j<T_{j,\max}$.
\end{itemize}

\subsubsection{Induktive Last}

\begin{itemize}
\item Wenn Schalter aus und kein Freilauf $\Rightarrow$ Überspannung.
\item Daher: Freilaufdiode oder antiparallele Diode ist Pflicht.
\end{itemize}

% ---------------------------------------------------------
\subsection{Gleichrichter (M1, B2, B6)}

\subsubsection{Systematik}

\begin{itemize}
\item M: Mittelpunkt (2 Dioden/Thyristoren)
\item B: Brücke (4 bzw. 6 Ventile)
\item Zahl: Pulszahl (1,2,6)
\item U: ungesteuert (Dioden)
\item C: gesteuert (Thyristoren)
\end{itemize}

\subsubsection{Einfluss von Zündwinkel $\alpha$}

\begin{itemize}
\item Wenn $\alpha\uparrow$ $\Rightarrow$ $\overline{u_d}\downarrow$.
\item B2C und B6C: $\overline{u_d}\propto \cos\alpha$.
\item Wenn $\alpha>90^\circ$ $\Rightarrow$ $\cos\alpha<0$ $\Rightarrow$ Rückspeisebetrieb möglich (bei passender Last).
\end{itemize}

\subsubsection{C-Glättung vs L-Glättung}

\begin{itemize}
\item C-Glättung: Spannung glatt, Strom impulsförmig, Spitzenstrom hoch, Netzrückwirkung schlecht.
\item L-Glättung: Strom glatt, Spannung wellig, Netzrückwirkung besser, Industrie typisch.
\item Wenn $C\uparrow$ $\Rightarrow$ $\Delta u \downarrow$ aber Spitzenstrom $\uparrow$.
\item Wenn $L\uparrow$ $\Rightarrow$ $\Delta i \downarrow$ und Strom wird kontinuierlicher.
\end{itemize}

\subsubsection{Kommutierungsüberlappung}

\begin{itemize}
\item Wenn Netzinduktivität $L_\sigma\uparrow$ oder $i_d\uparrow$ $\Rightarrow$ Überlappung $\mu\uparrow$.
\item $\mu\uparrow \Rightarrow \overline{u_d}\downarrow$.
\end{itemize}

\subsubsection{Harmonische}

\begin{itemize}
\item B6: typische Netzstromharmonische $n=6k\pm 1$.
\item B2: harmonisch stärker belastend als B6.
\end{itemize}

% ---------------------------------------------------------
\subsection{Wechselstromsteller (Phasenanschnitt)}

\begin{itemize}
\item Wenn $\alpha\uparrow$ $\Rightarrow$ Leitintervall kürzer $\Rightarrow$ $U_{\mathrm{RMS}}\downarrow$.
\item $\overline{u}$ bleibt bei symmetrischem AC-Anschnitt $=0$.
\item RMS sinkt nicht linear, sondern nach der bekannten Integralform.
\end{itemize}

\begin{itemize}
\item Wechselstrom\textbf{schalter}: nur EIN/AUS ganze Halbwellen.
\item Wechselstrom\textbf{steller}: Phasenanschnitt, Teilhalbwellen.
\end{itemize}

% ---------------------------------------------------------
\subsection{Gleichstromsteller (Chopper)}

\subsubsection{1-Quadrant (Buck-artig)}

\begin{itemize}
\item Wenn Tastgrad $d\uparrow$ $\Rightarrow$ $\overline{u_d}\uparrow$.
\item $\overline{u_d}=dU_{DC}$.
\item Wenn $L$ gross $\Rightarrow$ Strom nahezu konstant.
\end{itemize}

\subsubsection{2- und 4-Quadrant}

\begin{itemize}
\item 2Q: Energie kann in beide Richtungen, Spannung aber nur positiv (typisch Bremsbetrieb möglich).
\item 4Q: Spannung und Strom können Vorzeichen wechseln (Motor + Generator in beide Drehrichtungen).
\end{itemize}

% ---------------------------------------------------------
\subsection{DC-DC-Wandler (Buck, Boost, Buck-Boost)}

\subsubsection{Buck}

\begin{itemize}
\item $d\uparrow \Rightarrow u_2\uparrow$.
\item $0<d<1 \Rightarrow u_2 < U_{DC}$.
\item $L\uparrow \Rightarrow \Delta i_L \downarrow$.
\end{itemize}

\subsubsection{Boost}

\begin{itemize}
\item $d\uparrow \Rightarrow u_2\uparrow$ stark.
\item $d\rightarrow 1 \Rightarrow u_2 \rightarrow \infty$ (ideal), real begrenzt durch Verluste.
\end{itemize}

\subsubsection{Betriebsart}

\begin{itemize}
\item Wenn $L < L_{\min}$ $\Rightarrow$ diskontinuierlicher Betrieb $\Rightarrow$ Mittelwertformeln ändern sich.
\end{itemize}

% ---------------------------------------------------------
\subsection{Einphasiger Wechselrichter}

\begin{itemize}
\item Rechteckbetrieb: einfach, aber viele Harmonische.
\item PWM: Grundschwingung steuerbar, Harmonische hochfrequent.
\item Wenn Modulationsindex $m_a\uparrow$ $\Rightarrow$ Grundschwingung $\uparrow$ (linear bis $m_a\le 1$).
\end{itemize}

% ---------------------------------------------------------
\subsection{Dreiphasiger Wechselrichter}

\begin{itemize}
\item Dreiphasig: Leistungsfluss nahezu konstant.
\item 180$^\circ$-Betrieb: jede Phase 180$^\circ$ leitend.
\item 120$^\circ$-Betrieb: jede Phase 120$^\circ$ leitend, oft geringere Verluste.
\item Leiterspannung ist Differenz zweier Phasenspannungen: $u_{LL}=u_a-u_b$.
\end{itemize}

% ---------------------------------------------------------
\subsection{PWM (Sinus-PWM) – Spektrum und Auswirkungen}

\begin{itemize}
\item $f_s\uparrow \Rightarrow$ Ripplefrequenz $\uparrow$, Filterung einfacher, aber Schaltverluste $\uparrow$.
\item Spektrum: Seitenbänder um Vielfache von $f_s$.
\item R--L Last wirkt als Tiefpass $\Rightarrow$ Strom wird sinusähnlicher als Spannung.
\end{itemize}

% ---------------------------------------------------------
\subsection{Direktumrichter}

\begin{itemize}
\item Kein Zwischenkreis: AC $\rightarrow$ AC direkt.
\item Wenn Kommutierung falsch $\Rightarrow$ Kurzschluss von Netzphasen oder Unterbrechung der Last.
\item Ausgangsfrequenz begrenzt: typischerweise $f_2 < \frac{1}{2}f_1$.
\item Leistungsfluss kann positiv oder negativ sein $\Rightarrow$ Rückspeisung möglich.
\end{itemize}

% ---------------------------------------------------------
\subsection{Gleichstrommaschine (Motorbetrieb, Anlauf, Feldschwächung)}

\subsubsection{Kernaussagen}

\[
u_a = R_a i_a + k\Phi \omega,\qquad M=k\Phi i_a
\]

\subsubsection{Ankerspannung}

\begin{itemize}
\item $u_a\uparrow \Rightarrow \omega\uparrow$.
\item $u_a\downarrow \Rightarrow \omega\downarrow$.
\end{itemize}

\subsubsection{Lastmoment}

\begin{itemize}
\item $M_L\uparrow \Rightarrow i_a\uparrow \Rightarrow \omega\downarrow$.
\end{itemize}

\subsubsection{Feld}

\begin{itemize}
\item $\Phi\downarrow$ (Feldschwächung) $\Rightarrow \omega\uparrow$ aber $M_{\max}\downarrow$.
\item $\Phi\uparrow \Rightarrow \omega\downarrow$.
\end{itemize}

\subsubsection{Anlauf}

\[
\omega=0 \Rightarrow E=0 \Rightarrow i_{a,0}=\frac{u_a}{R_a}
\]

\begin{itemize}
\item Anlaufstrom sehr hoch $\Rightarrow$ Begrenzung durch Stromregelung/Anlaufwiderstand.
\item Anlaufmoment: $M_0=k\Phi i_{a,0}$.
\end{itemize}

% ---------------------------------------------------------
\subsection{Harmonische und Spektralanalyse}

\begin{itemize}
\item Nichtsinusförmig $\Rightarrow$ Grundschwingung + Harmonische.
\item THD hoch $\Rightarrow$ Verzerrung hoch.
\item R--L Last dämpft höhere Harmonische stärker: $I_n=\frac{U_n}{\sqrt{R^2+(n\omega L)^2}}$.
\end{itemize}

% ---------------------------------------------------------
\subsection{Fourier (schnelle Auswahlregeln)}

\begin{itemize}
\item Gerade Funktion $\Rightarrow b_n=0$.
\item Ungerade Funktion $\Rightarrow a_n=0$.
\item Zeitverschiebung $\Rightarrow$ Phasenfaktoren $e^{jk\omega t_0}$.
\item Sprungstellen $\Rightarrow$ Gibbs, Dirichlet: Mittelwert der Sprungwerte.
\end{itemize}

% =========================================================
